\documentclass[11pt]{article}

\usepackage[margin=1in]{geometry}
\usepackage{amsmath, amssymb}
\usepackage{graphicx}
\usepackage{booktabs}
\usepackage{hyperref}
\usepackage{microtype}
\usepackage{xcolor}
\usepackage{titlesec}
\usepackage{tocloft}
\usepackage{parskip}
\usepackage{setspace}
\usepackage{float}

\setstretch{1.15}
\definecolor{linkblue}{HTML}{1E5AA6}
\hypersetup{
    colorlinks=true,
    linkcolor=linkblue,
    urlcolor=linkblue,
    citecolor=linkblue,
    pdfauthor={Jonathan Davanzo},
    pdftitle={EEG-Based BCI for Robotic Hand Control: Full Report}
}
\titleformat{\section}{\large\bfseries}{\thesection}{1em}{}
\titleformat{\subsection}{\normalsize\bfseries}{\thesubsection}{1em}{}
\titleformat{\subsubsection}{\normalsize\bfseries}{\thesubsubsection}{1em}{}
\titlespacing*{\section}{0pt}{1.2em}{0.6em}
\titlespacing*{\subsection}{0pt}{0.9em}{0.4em}
\titlespacing*{\subsubsection}{0pt}{0.8em}{0.3em}
\setlength{\cftbeforesecskip}{0.2em}

\title{\textbf{EEG-Based BCI for Robotic Hand Control: Full Report}}
\author{Jonathan Davanzo}
\date{}

\begin{document}
\maketitle
\tableofcontents
\newpage

\begin{abstract}
This report presents a real-time, uncertainty-aware brain--computer interface
(BCI) for EEG-based finger and hand action classification with safe robotic
actuation. We describe the end-to-end pipeline from data acquisition and
labeling through model training, evaluation, and live deployment, and outline
the safety and reproducibility mechanisms that enable reliable demonstrations.
\end{abstract}

\section{Introduction}
EEG-based BCIs enable hands-free interaction by decoding neural activity into
actions. However, noise, artifacts, and uncertainty can degrade reliability,
especially in live robotics. This project targets a lightweight, real-time
pipeline that provides calibrated confidence and uncertainty estimates to gate
actuation safely.

\section{System Overview}

This project implements a real-time, uncertainty-aware brain--computer interface (BCI) for EEG-based finger and hand action classification, designed for safe robotic actuation and interactive demonstration. The system processes raw EEG data acquired from a Muse 2 headset, performs online signal cleaning and feature extraction, trains and evaluates neural network models, and deploys real-time inference with calibrated confidence and Bayesian uncertainty estimates.

The complete pipeline is modular and consists of five primary stages:
\begin{enumerate}
    \item EEG streaming and recording
    \item Event annotation and window extraction
    \item Neural network training
    \item Model evaluation, calibration, and validation
    \item Real-time inference with safety gating
\end{enumerate}

All stages are orchestrated via an automated experiment runner while preserving reproducibility, subject anonymity, and traceable experiment metadata.

\section{Data Acquisition}

\subsection{EEG Hardware and Sampling}

EEG data are acquired using a Muse 2 EEG headset, which provides four channels: TP9, AF7, AF8, and TP10. Data are streamed at a nominal sampling rate of 256~Hz using the Lab Streaming Layer (LSL) protocol. The system automatically resolves and connects to the active EEG stream.

\subsection{Real-Time Streaming}

Incoming EEG samples are buffered using a sliding window approach. Each window spans 0.25~seconds (64 samples per channel). This window duration balances temporal resolution with computational latency, enabling real-time inference while preserving discriminative neural features.

\section{Signal Preprocessing}

\subsection{Temporal Compression}

To reduce redundant computation and stabilize downstream ICA fitting, temporal compression is applied by removing consecutive duplicate samples across channels. This preserves signal transitions while reducing effective sample counts.

\subsection{Independent Component Analysis}

Artifact removal is performed using FastICA applied to each compressed window:
\begin{enumerate}
    \item EEG samples are standardized using a z-score transformation.
    \item ICA is fitted once per session and reused thereafter.
    \item Independent components are analyzed using Welch power spectral density.
    \item Components with dominant low-frequency (	extless{}4~Hz) or high-frequency (	extgreater{}35~Hz) power over the 8--30~Hz band are treated as artifacts and suppressed.
    \item Cleaned EEG is reconstructed via inverse ICA transformation.
\end{enumerate}

This adaptive approach enables ocular and muscular artifact suppression without predefined templates.

\section{Feature Extraction}

From each cleaned EEG window, a full time-series segment is retained (64 timesteps $\times$ 4 channels). A cleaned sample stream is saved to 	exttt{eeg\_features.csv} for event review, and a simple channel-wise mean feature vector is computed for diagnostics:
\[
\mathbf{x} = [\mu_{TP9}, \mu_{AF7}, \mu_{AF8}, \mu_{TP10}]
\]

These features are computationally lightweight, suitable for real-time inference, and consistent across training and evaluation.

\section{Event Annotation and Labeling}

\subsection{Event-Marking Interface}

A dedicated event-marking user interface allows post-hoc labeling of EEG data. The interface supports click-and-drag duration marking, keyboard-based labeling, undo/redo functionality, autosave, and export to EEGLAB-compatible \texttt{.set} files.

\subsection{Label Schema}

To align with BCI competition standards and multi-task learning, labels are structured hierarchically.

\subsubsection{Finger Labels}
\begin{itemize}
    \item Thumb
    \item Index
    \item Middle
    \item Ring
    \item Pinky
    \item Rest (optional)
\end{itemize}

\subsubsection{Action Labels}
\begin{itemize}
    \item Open
    \item Close
    \item Rest / No movement
\end{itemize}

Each EEG window is assigned both a finger identity and an action label, enabling multi-head neural network training.

\section{Window Extraction}

Annotated EEG recordings are segmented into fixed-length labeled windows and stored in \texttt{eeg\_windows.npz} with shape $[N, T, C]$. A compact summary is saved to \texttt{eeg\_\allowbreak windows.csv}. Windows inherit labels from overlapping annotated events. Ambiguous windows are either resolved via majority labeling or excluded.

This windowed dataset serves as the sole input for model training, evaluation, and validation.

\section{Neural Network Architecture}

\subsection{Model Design}

The primary model, \texttt{CNNLSTMFingerActionNet}, is a multi-head neural network comprising:
\begin{itemize}
    \item A shared feature encoder
    \item A finger classification head
    \item An action classification head
\end{itemize}

Dropout layers are included throughout the network to support Bayesian uncertainty estimation via Monte Carlo dropout.

\subsection{Training Objective}

The model is trained using a joint loss:
\[
\mathcal{L} = \mathcal{L}_{\text{finger}} + \mathcal{L}_{\text{action}}
\]
where each term is a categorical cross-entropy loss.

\section{Training Procedure}

Data are split into training (80\%) and test (20\%) sets using subject-wise splits when available; otherwise stratified sampling is used. Per-channel normalization statistics are fitted exclusively on the training set to prevent data leakage. The model is trained using the Adam optimizer for a fixed number of epochs with dropout enabled.

All artifacts, including trained weights, scalers, and predictions, are saved to disk.

\section{Evaluation and Calibration}

\subsection{Deterministic Evaluation}

For standard metrics, dropout is disabled. Accuracy, confusion matrices, and Expected Calibration Error (ECE) are computed on the held-out test set.

\subsection{Bayesian Uncertainty Estimation}

Uncertainty is estimated using Monte Carlo dropout:
\begin{enumerate}
    \item Dropout remains active during inference.
    \item Multiple stochastic forward passes are performed.
    \item Predictive mean and standard deviation are computed.
\end{enumerate}

Confidence is defined as the maximum predictive mean probability, while uncertainty is defined as the mean predictive standard deviation.

\section{Data and Model Validation}

The Deepchecks tabular validation framework is used to detect data leakage, distribution shifts, label inconsistencies, and anomalous model behavior. Validation operates on the same windowed dataset used for training.

\section{Real-Time Inference and Safety Gating}

\subsection{Confidence-Weighted Actuation}

Live actuation is gated by an adaptive confidence threshold that increases with uncertainty. Actuation is permitted only when predictions are stable across consecutive frames.

\subsection{Velocity Scaling}

Motor velocity is scaled according to:
\[
v = \text{confidence} \times (1 - \text{uncertainty})
\]
ensuring conservative behavior under uncertainty.

\section{Experiment Logging and Reproducibility}

Experiments are logged using anonymous subject identifiers, deterministic experiment hashes, and step-wise metadata records. This enables full traceability while preserving subject privacy.

\section{Automation and Execution}

All pipeline stages are orchestrated using an automated experiment runner that verifies dependencies, executes steps in order, and generates evaluation reports and figures.

\section{Ethical and Safety Considerations}

Subject anonymity is strictly preserved. Robotic actuation defaults to a disabled state under uncertainty, and conservative thresholds are enforced during live demonstrations.

\section{Intended Applications}

The system is designed for science fairs, research demonstrations, grant submissions, educational BCI prototypes, and safe human--robot interaction research.

\section{Results}
This section summarizes model accuracy, calibration, and live performance
metrics. Replace the placeholders below with actual results once finalized.
\begin{itemize}
    \item Finger classification accuracy: \textit{TBD}
    \item Action classification accuracy: \textit{TBD}
    \item Calibration (ECE): \textit{TBD}
    \item Live inference latency: \textit{TBD}
\end{itemize}

\section{Discussion}
The modular pipeline supports rapid iteration and safe demonstrations, but
performance depends on data quality, labeling fidelity, and subject variability.
Future work includes richer features, subject-adaptive calibration, and
additional uncertainty-aware gating strategies.

\section{Conclusion}
This report documents a complete EEG-to-robotics pipeline that prioritizes
real-time performance, calibrated uncertainty, and reproducible experimentation.

\section*{Acknowledgments}
Thanks to mentors, collaborators, and participants who supported data
collection and feedback.

\begin{thebibliography}{9}
\bibitem{muse2}
Muse 2 Headset, InteraXon, Inc. Product documentation and specifications.
\bibitem{lsl}
Lab Streaming Layer (LSL). https://labstreaminglayer.readthedocs.io/
\bibitem{deepchecks}
Deepchecks: A validation framework for tabular data and models.
\end{thebibliography}

\appendix
\section{Appendix: Reproducibility Checklist}
\begin{itemize}
    \item Fixed train/test split with stratification
    \item Normalization fit on training set only
    \item Saved scalers, weights, and predictions
    \item Logged experiment metadata and hashes
\end{itemize}

\end{document}
